\documentclass[11pt]{article}
\usepackage[T1]{fontenc}
\usepackage[margin=1in]{geometry}
\usepackage{amsmath,amssymb,amsthm,microtype}
\usepackage[colorlinks=true,urlcolor=blue,citecolor=blue]{hyperref}

\newtheorem{theorem}{Theorem}
\newtheorem{lemma}{Lemma}
\newtheorem{remark}{Remark}
\newcommand{\FBL}{\operatorname{FBL}}

\title{Near-isometric chain lifting in reflexive, dual, and bidual Banach lattices}
\author{A partial result for arXiv:2103.08170}
\date{11 August 2026}

\begin{document}
\maketitle

\begin{abstract}
The open optimal-constant problem for projective free Banach lattices over
lattices remains unresolved in the unrestricted category.  We prove three
sharp relative results.  For every linearly ordered set \(L\),
\(\FBL\langle L\rangle\) has \(1^+\)-controlled lifts through quotients of
reflexive Banach lattices and through weak-star normal quotients of dual
Banach lattices.  Every quotient of an arbitrary Banach lattice admits the
same lift after passing to the bidual.  If \(L\) is countable, quotients of
\(C(K)\)-spaces admit an exact norm-preserving lift.  The proofs isolate the
remaining obstruction: finite \(1^+\)-lifts glue by compactness, but a
general Banach lattice supplies neither compact bounded fibers nor a
lattice-preserving descent from its bidual.
\end{abstract}

\section{Problem and notation}

For a lattice \(L\), write \(P_L=\FBL\langle L\rangle\), and denote its
canonical generators by \(\delta_\ell\), \(\ell\in L\).  Its universal
property says that every bounded lattice homomorphism \(s:L\to X\) extends
uniquely to a Banach-lattice homomorphism
\(\widehat s:P_L\to X\), with
\[
 \|\widehat s\|=\sup_{\ell\in L}\|s(\ell)\|.
\]
For linearly ordered \(L\), the source paper proves that \(P_L\) is
projective precisely for the countable bounded chains, but obtains only a
\(2^+\) upper bound and asks
for the optimal constants \cite{AMRR}.  Finite lattices are known to
generate \(1^+\)-projective free Banach lattices \cite{AR}.

For a closed ideal \(J\) in a Banach lattice \(X\), let
\(Q:X\to X/J\) be the quotient map.  We write
\(\kappa_X:X\to X^{**}\) for the canonical embedding.

\section{Compact gluing}

\begin{lemma}[finite-lift compactness]\label{lem:compact}
Let \(L\) be a linearly ordered set, let \(Z\) be a Banach lattice whose
closed ball \(RB_Z\) is compact for a Hausdorff topology \(\tau\), and
suppose the positive cone of \(Z\) is \(\tau\)-closed.  Let
\(q:Z\to Y\) be a lattice quotient such that every fiber
\[
 K_\ell=\{z\in RB_Z:qz=t_\ell\}\qquad(\ell\in L)
\]
is nonempty and \(\tau\)-closed.  If every finite subchain \(F\subset L\)
has representatives \(z_\ell\in K_\ell\), \(\ell\in F\), satisfying
\(z_\ell\le z_m\) whenever \(\ell\le m\), then the entire family
\((t_\ell)_{\ell\in L}\) has such representatives in \(RB_Z\).
\end{lemma}

\begin{proof}
Each \(K_\ell\) is compact, so \(\prod_{\ell\in L}K_\ell\) is compact.
For \(\ell\le m\), the condition \(z_\ell\le z_m\) is closed because the
positive cone is closed.  These closed conditions have the finite
intersection property by hypothesis.  Their total intersection is
therefore nonempty.
\end{proof}

\begin{theorem}[relative \(1^+\)-projectivity and bidual lifting]
\label{thm:main}
Let \(L\) be any linearly ordered set, \(P_L=\FBL\langle L\rangle\), and
\(\varepsilon>0\).
\begin{enumerate}
 \item If \(X\) is reflexive, then every lattice homomorphism
 \[
 T:P_L\longrightarrow X/J
 \]
 has a lift \(\widetilde T:P_L\to X\) satisfying
 \[
 Q\widetilde T=T,\qquad
 \|\widetilde T\|\le(1+\varepsilon)\|T\|.
 \]
 \item The same conclusion holds when \(X=E^*\) is a dual Banach lattice
 with its canonical weak-star closed cone and \(J\) is a weak-star closed
 ideal.
 \item For arbitrary \(X\) and \(J\), there is a bidual lift
 \(\widetilde T:P_L\to X^{**}\) such that
 \[
 Q^{**}\widetilde T=\kappa_{X/J}T,\qquad
 \|\widetilde T\|\le(1+\varepsilon)\|T\|.
 \]
\end{enumerate}
Consequently the infimal lifting constant in each of the first two
relative categories, and in the bidual sense of the third, is exactly
\(1\).
\end{theorem}

\begin{proof}
Put \(C=\|T\|\); the case \(C=0\) is immediate.  For each \(\ell\in L\),
set \(t_\ell=T\delta_\ell\).  If \(F\subset L\) is finite, then \(F\) is a
finite lattice.  The map \(t|_F:F\to X/J\) induces a lattice homomorphism
\[
 T_F:P_F\longrightarrow X/J,\qquad
 \|T_F\|=\max_{\ell\in F}\|t_\ell\|\le C.
\]
The \(1^+\)-projectivity theorem for finite lattices gives a lift of
\(T_F\) of norm at most \((1+\varepsilon)C\).  Its values on the
generators are representatives of the \(t_\ell\)'s, are ordered along
\(F\), and lie in the ball of radius
\(R=(1+\varepsilon)C\).

If \(X\) is reflexive, apply Lemma~\ref{lem:compact} to the weak topology
on \(RB_X\).  Quotient fibers are weakly closed, and the positive cone is
weakly closed.  We obtain an order-preserving map \(s:L\to RB_X\) with
\(Qs(\ell)=t_\ell\).  Because \(L\) is a chain, order preservation is
equivalent to preservation of its finite meets and joins.  The universal
property gives \(\widetilde T=\widehat s\), with
\(\|\widetilde T\|\le R\), and equality on generators gives
\(Q\widetilde T=T\).

For \(X=E^*\), use the weak-star topology.  Its norm ball is weak-star
compact and its positive cone is weak-star closed.  A weak-star closed
ideal \(J\) makes \(X/J\) a dual quotient and \(Q\) weak-star continuous,
so its fibers are weak-star closed.  The same argument applies.

For arbitrary \(X\), work in \(X^{**}\) with its weak-star topology and
replace each finite lift \(x_\ell\in X\) by
\(\kappa_Xx_\ell\).  Define
\[
 K_\ell=\{x^{**}\in RB_{X^{**}}:
 Q^{**}x^{**}=\kappa_{X/J}t_\ell\}.
\]
These are nonempty weak-star compact sets; the cone of \(X^{**}\) is
weak-star closed, and \(Q^{**}\) is weak-star continuous.  Compact
gluing produces an ordered \(s:L\to X^{**}\), and the universal property
gives the asserted bidual lift.

Finally, every lift has norm at least \(\|T\|\), since the quotient and
canonical embedding are contractive.  Allowing arbitrary
\(\varepsilon>0\) proves the last assertion.
\end{proof}

\section{Exact lifting through \texorpdfstring{\(C(K)\)}{C(K)}-quotients}

The compactness proof uses an arbitrarily small norm loss inherited from
the finite-lattice theorem.  For countable chains and AM-space quotients,
one can remove even that loss.

\begin{lemma}[relative ordered Tietze extension]\label{lem:tietze}
Let \(F\) be closed in a compact Hausdorff space \(K\).  If
\(a,b\in C(K)\), \(a\le b\), and \(f\in C(F)\) satisfies
\(a|_F\le f\le b|_F\), then there exists \(g\in C(K)\) such that
\[
 g|_F=f,\qquad a\le g\le b.
\]
\end{lemma}

\begin{proof}
By the ordinary Tietze theorem, choose \(h\in C(K)\) with \(h|_F=f\).
Then
\[
 g=(h\vee a)\wedge b
\]
is continuous and lies between \(a\) and \(b\).  On \(F\), the assumed
inequalities give \((f\vee a|_F)\wedge b|_F=f\), so \(g|_F=f\).
\end{proof}

\begin{theorem}[norm-one \(C(K)\)-quotient lifting]\label{thm:ck}
Let \(L\) be a countable linearly ordered set, \(F\) a closed subset of a
compact Hausdorff space \(K\), and \(R_F:C(K)\to C(F)\) the restriction
quotient.  Every lattice homomorphism
\[
 T:P_L\longrightarrow C(F)
\]
has a lattice-homomorphic lift \(\widetilde T:P_L\to C(K)\) with
\[
 R_F\widetilde T=T,\qquad \|\widetilde T\|=\|T\|.
\]
\end{theorem}

\begin{proof}
Set \(C=\|T\|\), \(t_\ell=T\delta_\ell\), and enumerate
\(L=\{\ell_1,\ell_2,\ldots\}\).  We recursively construct
\(g_n\in C(K)\) such that
\[
 g_n|_F=t_{\ell_n},\qquad \|g_n\|_\infty\le C,
\]
and \(g_i\le g_j\) whenever \(i,j\le n\) and
\(\ell_i\le\ell_j\).

Suppose \(g_1,\ldots,g_{n-1}\) are defined.  Let
\[
 a_n=\max\bigl(\{-C\mathbf1\}\cup
   \{g_j:j<n,\ \ell_j<\ell_n\}\bigr),
\quad
 b_n=\min\bigl(\{C\mathbf1\}\cup
   \{g_j:j<n,\ \ell_n<\ell_j\}\bigr).
\]
The previous order relations give \(a_n\le b_n\).  On \(F\), the fact
that \(\ell\mapsto t_\ell\) is order preserving gives
\[
 a_n|_F\le t_{\ell_n}\le b_n|_F.
\]
Lemma~\ref{lem:tietze} supplies \(g_n\) between these two bounds and
extending \(t_{\ell_n}\).  This completes the recursion.

The map \(s(\ell_n)=g_n\) is a bounded lattice homomorphism from the
chain \(L\) to \(C(K)\), with rank at most \(C\).  Its universal extension
\(\widetilde T=\widehat s\) has norm at most \(C\) and restricts to \(T\).
The reverse inequality follows from contractivity of \(R_F\).
\end{proof}

\section{Why this does not settle the unrestricted problem}

Theorem~\ref{thm:main}(3) shows that the finite-lattice estimates already
contain a coherent global lift; it naturally lives in \(X^{**}\).  A
general principle of local reflexivity cannot finish the proof: it is
linear and finite-dimensional, whereas the required descent must
simultaneously preserve the lattice order for all generators and commute
exactly with \(Q\).  Nor does diagonal convexification solve the problem.
Convex combinations preserve the order relations of a chain and the
quotient equations, but a bounded sequence of finite lifts need not have
a weakly convergent subnet in a nonreflexive space.

The recursive proof of Theorem~\ref{thm:ck} uses the special relative
insertion property of \(C(K)\).  For a general Banach-lattice quotient,
clipping a lift between ordered endpoint lifts gives the familiar factor
\(2\), since
\[
 \bigl\||a|\vee|b|\bigr\|\le\|a\|+\|b\|,
\]
and no near-isometric simultaneous insertion theorem is available.  Thus
the unrestricted question for countable bounded chains reduces sharply
to whether the bidual ordered lift can always be descended, with
arbitrarily small norm loss, through an arbitrary Banach-lattice quotient.

\begin{thebibliography}{9}
\bibitem{AMRR}
A.~Avil\'es, G.~Mart\'inez-Cervantes, J.~D.~Rodr\'iguez Abell\'an, and
A.~Rueda Zoca,
\emph{Free Banach lattices generated by a lattice and projectivity},
Proc. Amer. Math. Soc. 150 (2022), 2071--2082,
\url{https://doi.org/10.1090/proc/15802}.

\bibitem{AR}
A.~Avil\'es and J.~D.~Rodr\'iguez Abell\'an,
\emph{Projectivity of the free Banach lattice generated by a lattice},
Arch. Math. 113 (2019), 515--524,
\url{https://arxiv.org/abs/1903.01191}.
\end{thebibliography}

\end{document}
